% Introductory description of the comparison methods used in the anchored
% simulation outputs.  The RLP-suffix results use the anRLP implementation for
% RLP.  This file is intended to be included with \input{} or compiled as a
% standalone note.
\documentclass[11pt,a4paper]{article}
\usepackage[margin=0.75in]{geometry}
\usepackage{xcolor}

\begin{document}

\section{Anchored Comparison Methods}

All estimates in the anchored simulation output are reported on the
\(\log(\mathrm{RR})\) scale.  This note focuses on two implementation details:
which weights are used by each estimator, and which sources of uncertainty are
included in the reported standard error or confidence interval.  For the
RLP-suffix simulation results, \texttt{RLP} is based on the
\texttt{anRLP} implementation, where the robust log-Poisson variance is a joint
sandwich variance.

\subsection{Simulation Setting}

The RR simulations generate independent binary covariates
\(Z_1,\ldots,Z_q\), with \(q=3\) in the default runs.  In the IPD trial, each
covariate has marginal probability \(p_{\mathrm{IPD}}=0.5\).  The AD covariate
distribution is either fixed by a user-specified target mean or chosen through
the theoretical relative ESS equation so that the anchored weights have the
requested ESS ratio; when no ESS target is specified, the default AD covariate
probability is \(0.35\).  Treatment is randomized with \(\Pr(X=1)=0.5\).
Conditional event probabilities are generated from the RR parameterization used
in \texttt{getProbRR}: the log relative risk is constant within each trial,
\(\log RR_i=\alpha_g\), and the nuisance log odds-product model is
\(\beta_{0g}+\sum_{k=1}^q\beta_{kg}Z_{ik}\), with \(\beta_{kg}\) equal to the
scenario-specific slope divided by \(\sqrt q\).  The scenarios cross event
frequency and hypothesis status.  Under the null, \(\alpha_g=0\) for both IPD
and AD data.  Under the alternative, the IPD log RR is \(0.4\) in the
common-event setting and \(0.7\) in the rare-event setting, whereas the
AD-generating log RR is \(0.7\) and \(1.1\), respectively.  The remaining
intercept and slope parameters are set separately for common versus rare events
to control the baseline event rate and covariate prognostic strength.

\subsection{Anchoring Weights}

Let \(Y_i\) be the binary outcome, \(X_i\in\{0,1\}\) the treatment indicator,
and \(Z_i\) the baseline covariates in the IPD trial.  The AD trial defines the
target covariate distribution.  The code first estimates the AD target mean
\[
  \widehat\mu_{\mathrm{AD}}
  =
  \frac{N_{0,\mathrm{AD}}\bar Z_{0,\mathrm{AD}}
        +N_{1,\mathrm{AD}}\bar Z_{1,\mathrm{AD}}}
       {N_{0,\mathrm{AD}}+N_{1,\mathrm{AD}}}.
\]
Its sampling variance is estimated as
\[
  \widehat{\mathrm{Var}}(\widehat\mu_{\mathrm{AD}})
  =
  \widehat\Sigma_{\mathrm{AD}}/m_{\mathrm{AD}},
\]
where \(\widehat\Sigma_{\mathrm{AD}}\) is the empirical covariance matrix of the
AD covariates and \(m_{\mathrm{AD}}\) is the AD sample size used for the target
moment.

The anchored estimators use exponential tilting weights
\[
  w_i(\widehat\gamma)
  =
  \frac{\exp\{\widehat\gamma^\top(Z_i-\widehat\mu_{\mathrm{AD}})\}}
       {n^{-1}\sum_{j=1}^n
       \exp\{\widehat\gamma^\top(Z_j-\widehat\mu_{\mathrm{AD}})\}} .
\]
These weights are normalized to have sample mean 1.  The calibration parameter
\(\widehat\gamma\) is obtained by minimizing
\[
  Q(\gamma)
  =
  \log\sum_{i=1}^n
  \exp\{\gamma^\top(Z_i-\widehat\mu_{\mathrm{AD}})\}.
\]
Equivalently, the fitted weights solve the balancing equation
\[
  n^{-1}\sum_{i=1}^n
  w_i(\widehat\gamma)(Z_i-\widehat\mu_{\mathrm{AD}})=0.
\]

\subsection{Uncertainty Layers}

For clarity, the variance calculations can be described in three layers.

\begin{itemize}
\item Layer 1 is the uncertainty in the final effect estimator after treating
the chosen weights and AD target mean as fixed.  Examples are the GLM coefficient
variance, robust Poisson sandwich variance, or fixed-weight CMH influence
function variance.
\item Layer 2 is the uncertainty from estimating the calibration parameter
\(\widehat\gamma\) in the IPD sample.
\item Layer 3 is the uncertainty in the AD target mean
\(\widehat\mu_{\mathrm{AD}}\), using
\(\widehat{\mathrm{Var}}(\widehat\mu_{\mathrm{AD}})
=\widehat\Sigma_{\mathrm{AD}}/m_{\mathrm{AD}}\).
\end{itemize}

Thus, a ``fixed-weight'' variance includes only Layer 1.  A full anchored
sandwich variance includes Layers 1--3.

\subsection{BRM Model Used by \texttt{brm\_IPD} and \texttt{brm}}

The BRM estimators use the RR/log-odds-product parameterization.  Define
\[
  p_x(Z_i)=\Pr(Y_i=1\mid X_i=x,Z_i),\qquad x=0,1.
\]
The target model is
\[
  \log\frac{p_1(Z_i)}{p_0(Z_i)}
  =
  v_{a,i}^{\top}\alpha ,
\]
and the nuisance model is
\[
  \log\left\{
  \frac{p_0(Z_i)}{1-p_0(Z_i)}
  \frac{p_1(Z_i)}{1-p_1(Z_i)}
  \right\}
  =
  v_{b,i}^{\top}\beta .
\]
In the current RR simulations, \(v_a\) contains only an intercept, so
\(\alpha\) is the common \(\log(\mathrm{RR})\).  The nuisance design is
\(v_b=(1,Z^\top)^\top\).  Given weights \(w_i\), the weighted likelihood is
maximized over \((\alpha,\beta)\) by alternating optimization.

\subsection{Estimator-by-Estimator Summary}

\renewcommand{\arraystretch}{1.25}
\begin{center}
\footnotesize
\begin{tabular}{|p{0.13\textwidth}|p{0.43\textwidth}|p{0.34\textwidth}|}
\hline
Method & Point estimate & Variance / CI uncertainty layers \\
\hline
\texttt{brm\_IPD}
&
Unweighted brm in the IPD population, weight \(w_i=1\).
&
Layer 1 only.  The SE is the model-based inverse Fisher information. CI is the
normal Wald interval
\(\widehat\alpha\pm1.96\,\widehat{\mathrm{se}}\).
\\
\hline
\texttt{brm}
&
brm weighted by the anchored exponential tilting weights
\(w_i(\widehat\gamma)\).
&
Layers 1--3.  The joint sandwich includes the brm score, the calibration
equation for \(\widehat\gamma\), and uncertainty in
\(\widehat\mu_{\mathrm{AD}}\).  CI is a normal Wald interval based on this joint
sandwich SE.
\\
\hline
\texttt{LB}
&
Log-binomial GLM weighted by \(w_i(\widehat\gamma)\).  The estimate is the
coefficient of \(X\).
&
Layer 1 only.  The model-based GLM SE treats \(w_i(\widehat\gamma)\) and
\(\widehat\mu_{\mathrm{AD}}\) as fixed. CI is the normal Wald GLM interval.
\\
\hline
\texttt{LP}
&
Log-Poisson GLM weighted by \(w_i(\widehat\gamma)\).  The estimate is the
coefficient of \(X\).
&
Layer 1 only.  The Poisson model-based SE treats
\(w_i(\widehat\gamma)\) and \(\widehat\mu_{\mathrm{AD}}\) as fixed. CI is the
normal Wald GLM interval.
\\
\hline
\texttt{RLP}
&
Robust log-Poisson GLM weighted by \(w_i(\widehat\gamma)\).  The estimate is
the coefficient of \(X\).
&
Layers 1--3 for the RLP-suffix results.  The joint sandwich includes the robust
log-Poisson estimating equation, the calibration equation for
\(\widehat\gamma\), and uncertainty in \(\widehat\mu_{\mathrm{AD}}\).  CI is a
normal Wald interval.
\\
\hline
\texttt{CMH}
&
Stratified Mantel--Haenszel common log RR using weighted stratum counts:
\[
  \log
  \frac{\sum_s a_s n_{0s}/N_s}
       {\sum_s c_s n_{1s}/N_s},
\]
where the counts use \(w_i(\widehat\gamma)\).  Strata are observed covariate
combinations.
&
Layer 1 only.  The influence-function SE treats \(w_i(\widehat\gamma)\) as
fixed. CI is a normal Wald interval.
\\
\hline
\texttt{RCMH}
&
Same point estimate as \texttt{CMH}.
&
Layers 1--3.  The sandwich adds uncertainty from estimating
\(\widehat\gamma\) and \(\widehat\mu_{\mathrm{AD}}\) to the CMH influence
function. CI is a normal Wald interval.
\\
\hline
\texttt{brmad}
&
Boundary-adjusted \texttt{brm}.  It equals \texttt{brm} unless a weighted event
rate is 0 or 1; then \(\log(p_1/p_0)\) is estimated from Jeffreys beta
posteriors using weighted counts.
&
Non-boundary case: same Layers 1--3 Wald CI as \texttt{brm}.  Boundary case:
posterior SD and posterior quantile interval, treating weighted counts as fixed.
\\
\hline
\texttt{GC}
&
STC-type g-computation.  First fit a logistic outcome model in the IPD.  Then
sample \(Z_j^\star\) from the AD covariate data and predict risks under
\(X=1\) and \(X=0\).  No exponential tilting weights are used.
&
IPD bootstrap layer only.  In each bootstrap replicate, resample IPD rows, refit
the logistic model, and recompute \(\widehat\alpha_{\mathrm{GC}}\) on the same
fixed \(Z^\star\).  The reported SE is the sample SD of these bootstrap
estimates.  The CI is Wald type.
\textcolor{red}{We will add bootstrap resampling of the AD data.  The present
results bootstrap the IPD and sample \(Z_j^\star\) only once.}
\\
\hline
\texttt{brm\_bc}
&
Same point estimate as \texttt{brm}.  The CI is obtained by inverting a
weighted profile-LRT using \(w_i(\widehat\gamma)\).
&
The reported SE is copied from \texttt{brm}.  The profile-LRT CI itself is
conditional on fixed \(w_i(\widehat\gamma)\) and
\(\widehat\mu_{\mathrm{AD}}\).  The \(p\)-value uses the same profile-LRT test
at \(\alpha=0\).
\\
\hline
\texttt{brm\_ad\_bc}
&
Same point estimate as \texttt{brmad}.  The CI is obtained by the same weighted
profile-LRT inversion as \texttt{brm\_bc}.
&
The reported SE is copied from \texttt{brmad}.  The profile-LRT CI is
conditional on fixed \(w_i(\widehat\gamma)\) and
\(\widehat\mu_{\mathrm{AD}}\).  The \(p\)-value uses the same profile-LRT test
at \(\alpha=0\).
\\
\hline
\end{tabular}
\end{center}

\subsection{Details of \texttt{RLP} in the RLP-Suffix Results}

The RLP-suffix simulation results use the \texttt{anRLP} implementation of
\texttt{RLP}, through
\texttt{fit\_robust\_lp\_joint()}.  The fitted outcome model is the weighted
log-Poisson model
\[
  \log E(Y_i\mid X_i,Z_i)
  =
  X_i\beta_X+\beta_0+Z_i^\top\beta_Z ,
\]
using the anchored weights \(w_i(\widehat\gamma)\).  The point estimate reported
as \texttt{RLP} is \(\widehat\beta_X\).

The variance is a joint sandwich variance, not the usual fixed-weight HC0
variance.  Let \(M_i=(X_i,1,Z_i^\top)^\top\) denote the model matrix row,
\(\widehat\mu_i=\exp(M_i^\top\widehat\beta)\), and
\[
  \psi_{\beta i}
  =
  w_i(\widehat\gamma)M_i(Y_i-\widehat\mu_i)
\]
be the weighted Poisson estimating function.  For the calibration equation, let
\[
  D_i=Z_i-\widehat\mu_{\mathrm{AD}},\qquad
  H_i=D_i-\frac{\sum_j w_jD_j}{\sum_j w_j}.
\]
The code forms the calibration sensitivity
\[
  A_\gamma
  =
  n^{-1}\sum_i w_i D_iH_i^\top
\]
and the cross-sensitivity of the Poisson score to the calibration parameter,
\[
  A_{\beta\gamma}
  =
  n^{-1}\sum_i w_i(Y_i-\widehat\mu_i) M_iH_i^\top .
\]
These terms are used to modify the influence function for \(\widehat\beta\), so
the IPD part of the covariance, \(\widehat V_{\mathrm{IPD}}\), includes both the
usual robust Poisson uncertainty and the uncertainty from estimating
\(\widehat\gamma\).

When \texttt{account\_for\_target\_mean=TRUE}, as in the current anchored call,
the method also adds the AD target-mean component
\[
  \widehat V_{\mu}
  =
  \frac{\partial\widehat\beta}{\partial\widehat\mu_{\mathrm{AD}}}
  \widehat{\mathrm{Var}}(\widehat\mu_{\mathrm{AD}})
  \left(
  \frac{\partial\widehat\beta}{\partial\widehat\mu_{\mathrm{AD}}}
  \right)^\top .
\]
Therefore the reported RLP covariance is
\[
  \widehat V_{\mathrm{RLP}}
  =
  \widehat V_{\mathrm{IPD}}+\widehat V_{\mu},
\]
and the reported SE for \texttt{RLP} is the square root of the
\(\beta_X\) diagonal element of this matrix.  In the layer terminology above,
the \texttt{anRLP} version of \texttt{RLP} includes Layers 1--3.

\subsection{Exact CI Computation}

For \texttt{brm\_bc} and \texttt{brm\_ad\_bc}, the interval endpoints are
not Wald endpoints.  For a candidate value \(a\), the code fixes
\(\alpha=a\), profiles over the nuisance parameter \(\beta\), and computes the
weighted profile likelihood ratio statistic
\[
  T_{\mathrm{obs}}(a)
  =
  2\{\ell_w(\widehat\alpha,\widehat\beta)
  -
  \ell_w(\alpha=a,\widehat\beta_a)\}.
\]
It then simulates outcomes under the constrained fitted model and recomputes the
same statistic to obtain \(\{T_1(a),\ldots,T_M(a)\}\).  The acceptability of
\(a\) is a two-sided Monte Carlo tail-area measure.  Values with acceptability
at least 0.05 are included in the 95\% confidence set.  The implementation first
searches over \(\widehat\alpha\pm4\widehat{\mathrm{se}}\), truncated to
\([-8,8]\), and then refines the endpoints by bisection.

Importantly, this profile-LRT simulation is conditional on the fitted anchored
weights.  Therefore, although the displayed SE for \texttt{exact\_an} is copied
weights.  Therefore, although the displayed SE for \texttt{brm\_bc} is copied
from \texttt{brm}, the exact interval endpoints themselves do not integrate
over uncertainty in \(\widehat\gamma\) or \(\widehat\mu_{\mathrm{AD}}\).

\subsection{Method Names}

In this note, the display names correspond to the raw simulation columns as
follows:
\[
\begin{array}{llll}
\texttt{brm\_IPD}=\texttt{brm},&
\texttt{brm}=\texttt{brm\_an},&
\texttt{LB}=\texttt{LB\_an},&
\texttt{LP}=\texttt{LP\_an},\\
\texttt{RLP}=\texttt{RLP\_an},&
\texttt{CMH}=\texttt{CMH\_an},&
\texttt{RCMH}=\texttt{CMH\_sandwich\_an},&
\texttt{brmad}=\texttt{brmad\_an},\\
\texttt{GC}=\texttt{GC},&
\texttt{brm\_bc}=\texttt{exact\_an},&
\texttt{brm\_ad\_bc}=\texttt{ad\_exact\_an}.&
\end{array}
\]

\end{document}
